% ar/book.tex -- the master file for the Arabic (Modern Standard) edition.
% Hand-edited. Build with: finishing/tools/build_edition.sh ar
%
% Shared by every edition: the preamble (switched by \booklanguage,
% \booklocale and \bookmainfont), edition-common.tex (English captions, notes,
% heads) and edition-body.tex (the document). ar/sections.tex and
% ar/status.tex are GENERATED by finishing/tools/editions.py.
% Right to left: babel's bidi=basic sets the direction by language, so a Latin
% run (a citation, an English title) sits left to right inside the line.
% Arabic is set in Noto Naskh Arabic (a system font) through HarfBuzz. biblatex
% has no Arabic strings; the edition's own are in ar/biblatex-ar.lbx.
%
% The title, subtitle and notice are a language model's drafts, not yet
% reviewed by a native reader; the title itself is an open question (D-810).
\documentclass[11pt,oneside]{book}
\def\bookcode{ar}
\def\booklanguage{arabic}
\def\booklocale{arabic}
\PassOptionsToPackage{bidi=basic}{babel}
\input{ar/status}
\newcommand{\booklegal}{%
     حقوق النشر \textcopyright{} 2026 \bookauthor.\\
     منشور بموجب رخصة \foreignlanguage{english}{CC BY-NC-ND 4.0}.\\[0.4\baselineskip]
     تبقى مواد الأطراف الثالثة المقتبسة في هذا العمل، بما فيها العبارات الافتتاحية،\\
     ملكًا لأصحاب حقوقها، وتُنقل بموجب مبدأ الاستخدام العادل (fair use).\\
     ولا تسري الرخصة المذكورة أعلاه عليها.\\[0.4\baselineskip]
     الترجمة العربية جارية. تَرِد الأقسام التي لم تُترجم بعد بالإنجليزية.\\
     الأقسام المترجمة مسوّدات أعدّها نموذج لغوي،\\
     ولم يراجعها بعدُ متحدث أصلي باللغة.%
     \ifx\editionstalelist\empty\else\\[0.4\baselineskip]
     تَرِد هذه الأقسام بالإنجليزية إلى أن تُحدَّث ترجمتها:\\
     \editionstalelist.\fi}
\input{preamble}
% Noto's faces for this script carry no Latin letters, digits or ASCII
% punctuation; luaotfload's fallback takes any glyph the face lacks from the
% English book's face, so names, numbers and citations print.
\directlua{luaotfload.add_fallback("aflatin", {"TeX Gyre Pagella:mode=harf;"})}
\babelfont[arabic]{rm}[Renderer=HarfBuzz, RawFeature={fallback=aflatin}]{Noto Naskh Arabic}
\DeclareLanguageMapping{arabic}{ar/biblatex-ar}
% babel switches Latin runs, and the ASCII hyphen (which the face lacks and the
% fallback can't reach, since TeX handles an explicit hyphen as a hyphenation
% point), to English and the English face.
\babelprovide[onchar=ids fonts]{english}
\babelcharproperty{`-}{locale}{english}
% The same for the edition's own script: its characters in an English section
% (the bibliography strings in a citation) switch to its locale and face.
\babelprovide[onchar=ids fonts]{arabic}
\input{edition-common}

\newcommand{\booktitlemain}{ذكاء مناهض للفاشية}
\newcommand{\booksubtitle}{آلات قد ترفض، والجمهور الذي ينبغي أن يشكّلها}
\newcommand{\bookauthor}{Joshua G. Stern}
\newcommand{\bookreferences}{المراجع}
\title{\booktitlemain: \booksubtitle}
\author{\bookauthor}
\date{}
\hypersetup{pdftitle={\booktitlemain: \booksubtitle}, pdfauthor={\bookauthor}}
\input{edition-body}
